\documentclass{article}
\usepackage{amsmath,amsthm}
\newtheorem{definition}{Definition}
\newtheorem{proposition}{Proposition}
\title{The Roofline Model in One Page}
\begin{document}
\maketitle

\section{Two limits}
A kernel is limited either by arithmetic or by moving bytes.
The roofline model compares the two limits for one machine and one kernel.

\subsection{Machine parameters}
Let $P$ be the peak arithmetic throughput in FLOP/s and $B$ the memory bandwidth in bytes/s.

\subsection{Kernel parameters}
\begin{definition}
The \emph{arithmetic intensity} of a kernel is $I = F / Q$, where $F$ is the number of floating-point operations and $Q$ the number of bytes moved to and from memory.
\end{definition}

\section{The bound}
\begin{proposition}
The attainable throughput of a kernel is at most $\min(P, I \cdot B)$.
\end{proposition}
\begin{proof}
The kernel cannot exceed the peak $P$. Moving $Q$ bytes takes at least $Q / B$ seconds, so the $F$ operations run at most at $F / (Q/B) = I \cdot B$ operations per second.
\end{proof}

The two limits meet at the ridge point $I^\ast = P / B$. Kernels with $I < I^\ast$ are memory-bound and kernels with $I > I^\ast$ are compute-bound.

\section{Typical kernels}
\begin{table}[h]
\centering
\begin{tabular}{lr}
Kernel & Intensity (FLOP/byte) \\
\hline
Vector addition & 0.08 \\
Softmax & about 1 \\
Matrix multiply, large & hundreds \\
Single-token decode & about 1 \\
\end{tabular}
\caption{Approximate arithmetic intensity in 16-bit precision.}
\end{table}

\section{Prefill versus decode}
Prefill multiplies every loaded weight into many tokens, so its intensity is high and it is compute-bound.
Decode generates one token at a time, reads every weight once per token and is memory-bound, which is why smaller weights make decode faster.

\end{document}
